\chapter{Description}

\section{Figures}

\subsection{Basic Example}

\begin{figure}[H]
    \centering
    % \includegraphics[width=\textwidth]{figures/image}
    \rule{6cm}{6cm}
    \caption{Example Figure}
    \label{fig:error_function}
\end{figure}

Following shows the code to the basic example for including a Figure. The line that includes the actual is due to a lack of one outcommented and replaced by a black box.

\begin{verbatim}
\begin{figure}[H]
    \centering
    % \includegraphics[width=\textwidth]{figures/image}
    \rule{6cm}{6cm}
    \caption{Example Figure}
\label{fig:error_function}
\end{figure}

\end{verbatim}

\subsection{Required Packages}


\begin{verbatim}
\usepackage{graphicx}   % To include images
\usepackage{float}      % strict placement control using \texttt{[H]}
\end{verbatim}

\subsection{Scaling and Sizing}

You can control image size using:

\begin{verbatim}
\includegraphics[width=0.5\textwidth]{figures/image}
\includegraphics[height=5cm]{figures/image}
\end{verbatim}

\section{Tables}

\subsection{Basic Example}

\begin{table}[H]
    \centering
    \caption{Basic Table Example}
    \label{tab:basic}
    \begin{tabular}{|l|c|r|}
        \hline
        Left & Center & Right \\ \hline
        A & B & C \\ \hline
    \end{tabular}
\end{table}

The following shows the corresponding code:

\begin{verbatim}
\begin{table}[H]
    \centering
    \caption{Basic Table Example}
    \label{tab:basic}
    \begin{tabular}{|l|c|r|}
        \hline
        Left & Center & Right \\ \hline
        A & B & C \\ \hline
    \end{tabular}
\end{table}
\end{verbatim}

\subsection{Required Packages}

\begin{verbatim}
\usepackage{float}     % Enables [H] placement
\usepackage{makecell}  % Allows line breaks inside cells
\usepackage{graphicx}  % Required for \resizebox
\end{verbatim}

\subsection{Column Alignment}

The \texttt{tabular} column specification defines alignment:

\begin{center}
\begin{tabular}{ll}
\texttt{l} & Left-aligned\\
\texttt{c} & Centered\\
\texttt{r} & Right-aligned\\
\texttt{|} & Vertical line\\
\end{tabular}
\end{center}

Example:
\begin{verbatim}
\begin{tabular}{|l|c|r|}
\end{verbatim}

\subsection{Scaling and Width Control}

Wide tables can be scaled using:

\begin{verbatim}
\resizebox{\textwidth}{!}{%
\begin{tabular}{...}
...
\end{tabular}
}
\end{verbatim}

This forces the table to fit the page width.

\subsection{Line Breaks in Cells}

To include line breaks inside cells:

\begin{verbatim}
\makecell[l]{Line 1 \\ Line 2}
\end{verbatim}

This requires the \texttt{makecell} package.

\subsection{Large / Complex Table Example}

\begin{table}[H]
    \centering
    \caption{Example for big table}
    \label{tab:bigTable}
    \resizebox{\textwidth}{!}{
    \begin{tabular}{|l|l|l|l|l|l|}
        \hline
        \textbf{Source} & \textbf{Summary} & \textbf{Link to Source} & \textbf{Priority} & \textbf{Target Date} & \textbf{Status} \\ \hline
        Codebeamer & Summary in Tool, name in JSON & https://Codebeamer.azure.zf-world.com/item/\{itemID\} & priority & endDate & status \\ \hline
        MS Planner & Title & \makecell[l]{
        https://teams.microsoft.com/l/entity/com.microsoft.teamspace.tab.planner\\
        /mytasks?tenantId=eb70b763-b6d7-4486-8555-8831709a784e\&webUrl=https\%3A\\
        \%2F\%2Ftasks.teams.microsoft.com\%2Fteamsui\%2FpersonalApp\%2Falltaskli\\
        sts\&context=\{"subEntityId":"/v1/plan/\{planId\}/view/board/task/{id}"\}\\ \\ (replace \{planId\}, \{id\} with body.value.planId and body.value.id)} & priority & dueTime & percent-Complete \\ \hline
        Windchill & \makecell[l]{In Item-Field with \\ \texttt{Name} = \texttt{Summary}; summary stands \\ in \texttt{Item-Field.short-text.value}} & \makecell[l]{https://skobde-mks-im.kobde.trw.com:7001/im/viewissue?selection=\{itemID\} \\ integrity://skobde-mks-im.kobde.trw.com:7001/im/viewissue?selection=\{itemID\}} & ALM\_Priority & CompletionDate & State \\ \hline
    \end{tabular}}
\end{table}

\section{Float Placement Options for Figures and Tables}

Figures and Tables are floats, and their placement is controlled by optional parameters:

\begin{center}
\begin{tabular}{ll}
\texttt{h} & Place approximately here \\
\texttt{t} & Top of page \\
\texttt{b} & Bottom of page \\
\texttt{p} & Separate float page \\
\texttt{!} & Override internal constraints \\
\texttt{H} & Place exactly here (requires \texttt{float} package) \\
\end{tabular}
\end{center}

Example:
\begin{verbatim}
\begin{figure}[!htb]
\end{figure}

\begin{table}[!htb]
\end{table}
\end{verbatim}
\section{Referencing and Correct Citing}

\subsection{Referencing Figures and Tables}

\begin{verbatim}
As shown in Figure~\ref{fig:error_function}.
As listed in Table~\ref{tab:basic}.
See Section~\ref{sec:tables}.
\end{verbatim}

The tilde (\texttt{\textasciitilde}) creates a non-breaking space to prevent line breaks between the word and the reference number. Ensure that corresponding labels are defined.

\subsection{Citing Sources}

Citing is used to reference external literature and sources. \autocite{example}

\subsubsection{Using biblatex}

A more flexible modern approach:

\begin{verbatim}
\usepackage[backend=biber]{biblatex}
\addbibresource{references.bib}
\end{verbatim}

Cite sources in the text:

\begin{verbatim}
\autocite{key}
\end{verbatim}

Print the bibliography:

\begin{verbatim}
\printbibliography
\end{verbatim}

Or alternatively:

\begin{verbatim}
\printbibheading
\printbibliography[
    nottype=online,
    heading=subbibliography,
    title={Printed sources}
]
\printbibliography[
    type=online,
    heading=subbibliography,
    title={Online sources}
]
\end{verbatim}

\subsection{Example Bibliography Entry}

\begin{verbatim}
@article{example,
    author  = {Author, A.},
    title   = {Example Title},
    journal = {Journal Name},
    year    = {2024}
}
\end{verbatim}

\section{Acronyms and Glossaries}

Acronyms and glossary entries, such as \gls{ki} or \gls{cpu}, can be referenced in the text using the \texttt{gls} command:

\begin{verbatim}
\gls{cpu}
\end{verbatim}

Typically, acronyms and glossary entries are defined in a separate file (e.g., \texttt{acronyms.tex}), which is included in the main document using:

\begin{verbatim}
\newacronym{ki}{KI}{Künstliche Intelligenz}

\newglossaryentry{cpu}{
  name={CPU},
  description={Central Processing Unit}
}
\end{verbatim}

In this file, entries are defined as follows:

\begin{verbatim}
\newacronym{ki}{KI}{Künstliche Intelligenz}

\newglossaryentry{cpu}{
  name={CPU},
  description={Central Processing Unit}
}
\end{verbatim}

The glossary and acronym list can then be generated using:

\begin{verbatim}
\printglossaries
\end{verbatim}

When an acronym is used for the first time, the full term is automatically printed, followed by the acronym in parentheses. Subsequent uses will display only the acronym.

\section{To Add}

\begin{itemize}
    \item long tables
    \item create plots from cvs data
    \item pictures with tikz
    \item equations
    \item Code
\end{itemize}

\section{Planned Changes}

\begin{itemize}
    \item Examples shall be removed from 'description'
    \item 'description' shall be renamed to 'foundations'
    \item new 'description' shall be added
    \item new 'description' shall contain a collection of examples out of previous works
\end{itemize}